% TODO checklist:
% - [x] Inspected src/security/authorization.py (RBAC logic, scope matching, wildcard support)
% - [x] Inspected src/security/perm.py (Permission utilities, role inference)
% - [x] Inspected src/agent_policies/roles.yaml (Role definitions and scopes)
% - [x] Inspected src/mcp/policies.yaml (Tool-level authorization policies)
% - [x] Reviewed Q&A.md Q17 (RBAC implementation details)
% - [x] Reviewed README.md (Authorization features overview)

\section{Authorization}
\label{sec:authorization}

The Cineca Agentic Platform implements a comprehensive Role-Based Access Control (RBAC) system that governs access to API endpoints, MCP tools, graph database operations, and administrative functions. The authorization subsystem is designed with defense-in-depth principles, ensuring that multiple layers of access control are applied throughout the request lifecycle.

\subsection{RBAC Model}
\label{subsec:rbac-model}

The authorization model is built on the concept of \textit{roles} that expand to \textit{scopes}. This two-tier approach provides both administrative simplicity (assigning roles) and fine-grained control (scope-level permissions).

\paragraph{Terminology}
\begin{itemize}
    \item \textbf{Roles}: Coarse-grained labels assigned to users (e.g., ``admin'', ``researcher'', ``operator''). Roles are typically assigned at the Identity Provider level and arrive in tokens via the \texttt{roles} claim.
    \item \textbf{Scopes}: Fine-grained permission strings following a \texttt{resource:action} pattern (e.g., ``tools:invoke'', ``graph:query'', ``admin:tenants:write''). Scopes represent specific capabilities.
    \item \textbf{Permissions}: In Auth0-style tokens, explicit permission arrays that are treated equivalently to scopes.
\end{itemize}

\paragraph{Role Hierarchy}
The platform defines a hierarchy of roles with increasing privilege levels:

\begin{enumerate}
    \item \textbf{guest}: Minimal capabilities for trial or external users. Limited to schema discovery and basic output formatting.
    \item \textbf{researcher}: Scientists with read access to graphs, model inference, and visualization capabilities. Cannot modify data.
    \item \textbf{curator}: Data stewards with graph CRUD operations, quality checks, and archive management.
    \item \textbf{data\_engineer}: ETL-focused role with bulk operations, system backup, and database switching capabilities.
    \item \textbf{operator}: SRE/operations role for monitoring, health checks, and rate limit management. No data mutation.
    \item \textbf{auditor}: Read-only security and compliance oversight with strict safety enforcement.
    \item \textbf{admin}: Full system access including cross-tenant operations and all administrative functions.
\end{enumerate}

\subsection{RBAC Implementation Details}
\label{subsec:rbac-implementation}

The authorization implementation in \texttt{src/security/authorization.py} provides a flexible and extensible framework for access control decisions.

\subsubsection{Role-to-Scope Expansion}

Roles are expanded to scopes using a configurable policy mapping. The platform supports multiple policy sources with a defined precedence:

\begin{enumerate}
    \item \textbf{Built-in Defaults}: Conservative baseline permissions.
    \item \textbf{src/agent\_policies/roles.yaml}: Primary role policy file.
    \item \textbf{src/mcp/policies.yaml}: Tool-specific policies (merged).
\end{enumerate}

\begin{lstlisting}[language=Python,caption={Role-to-scope expansion logic}]
_DEFAULT_ROLE_SCOPES: dict[str, list[str]] = {
    "user": [
        "read",
        "agent.run",
        "tools.invoke",
        "models.complete",
        "system.health",
    ],
    "admin": [
        "*",  # Full access via wildcard
    ],
}

def _expand_roles_to_scopes(scopes_or_roles: Iterable[str]) -> list[str]:
    """Expand each token scope/role using the role->scopes mapping."""
    expanded: list[str] = []
    seen = set()
    for token in scopes_or_roles:
        if token in _ROLE_TO_SCOPES:
            for s in _ROLE_TO_SCOPES[token]:
                if s not in seen:
                    expanded.append(s)
                    seen.add(s)
        elif token not in seen:
            expanded.append(token)
            seen.add(token)
    return expanded
\end{lstlisting}

\subsubsection{Scope Format and Patterns}

Scopes follow a hierarchical naming convention that supports wildcards:

\begin{itemize}
    \item \textbf{Exact Match}: \texttt{tools:invoke} matches only that specific scope.
    \item \textbf{Resource Wildcard}: \texttt{tools:*} matches any action on the tools resource.
    \item \textbf{Super Wildcard}: \texttt{*} grants access to everything (reserved for admin role).
\end{itemize}

The wildcard matching implementation uses regex-based pattern evaluation:

\begin{lstlisting}[language=Python,caption={Wildcard scope matching}]
def _scope_satisfies(candidate: str, required: str) -> bool:
    """Return True if candidate scope covers required."""
    if candidate == "*":
        return True
    # Convert wildcard pattern to regex
    pattern = r"^" + re.escape(candidate).replace(r"\*", ".*") + r"$"
    return bool(re.fullmatch(pattern, required))
\end{lstlisting}

\subsubsection{Authorization Check Modes}

The platform supports two authorization modes:

\begin{enumerate}
    \item \textbf{any} (default): At least one required scope must be satisfied by the user's effective scopes. Suitable for endpoints with alternative access paths.
    \item \textbf{all}: All required scopes must be satisfied. Used when multiple distinct permissions are needed.
\end{enumerate}

\begin{lstlisting}[language=Python,caption={Scope checking with mode support}]
def check_scopes(
    user_scopes_or_roles: Iterable[str],
    required: str | Iterable[str],
    *,
    mode: str = "any",
) -> bool:
    """Evaluate authorization based on scope matching."""
    req = _as_list(required)
    eff = _expand_roles_to_scopes(user_scopes_or_roles)
    
    if "*" in eff:
        return True  # Super admin bypass
    
    if mode == "any":
        for r in req:
            for c in eff:
                if _scope_satisfies(c, r):
                    return True
        return False
    
    # mode == "all"
    return all(any(_scope_satisfies(c, r) for c in eff) for r in req)
\end{lstlisting}

\subsubsection{Authorization Decision with Audit Trail}

Every authorization decision is captured with an audit trail:

\begin{lstlisting}[language=Python,caption={Authorization with audit logging}]
def authorize(
    user: Any,
    required_scopes: str | Iterable[str],
    *,
    resource: str | None = None,
    action: str | None = None,
    mode: str = "any",
) -> AuthzDecision:
    """Make an authorization decision and emit an audit event."""
    req = tuple(_as_list(required_scopes))
    raw = _extract_scopes(user)
    eff = tuple(_expand_roles_to_scopes(raw))
    allowed = check_scopes(raw, req, mode=mode)
    principal = _principal(user)
    
    reason = (
        "granted by wildcard" if "*" in eff
        else ("all scopes satisfied" if allowed and mode == "all"
              else ("one scope satisfied" if allowed 
                    else "insufficient scopes"))
    )
    
    # Emit audit event
    audit_policy_decision(
        policy="role_scope_policy",
        subject=principal,
        action=action or req[0],
        resource=resource or "unknown",
        allowed=allowed,
        reason=reason,
        attributes={"required": list(req), "effective_scopes": list(eff)}
    )
    
    return AuthzDecision(
        allowed=allowed, 
        reason=reason, 
        required=req, 
        effective_scopes=eff
    )
\end{lstlisting}

\subsubsection{FastAPI Integration}

The authorization system integrates seamlessly with FastAPI through dependency injection:

\begin{lstlisting}[language=Python,caption={FastAPI authorization dependency}]
def require_scopes(
    required_scopes: str | Iterable[str],
    *,
    mode: str = "any",
    resource: str | None = None,
):
    """Create a dependency that enforces scope requirements."""
    async def _dep(user=Depends(get_current_user)):
        authorize_or_403(
            user,
            required_scopes,
            resource=resource or "route",
            mode=mode,
        )
        return user
    return _dep

# Usage in routers
@router.get("/admin")
async def admin_only(user=Depends(require_scopes("admin:all"))):
    return {"ok": True}
\end{lstlisting}

\subsubsection{Permission Utilities}

The \texttt{perm.py} module provides additional utilities for working with permissions:

\begin{lstlisting}[language=Python,caption={Permission extraction and checking}]
def current_permissions(user) -> set[str]:
    """Extract effective permissions from a Principal-like object.
    
    Precedence: permissions -> scope/scopes -> roles mapping.
    """
    perms: set[str] = set()
    raw = getattr(user, "raw", {}) or {}
    
    # 1) Explicit permissions claim (Auth0-style)
    if isinstance(raw.get("permissions"), (list, tuple)):
        perms.update(raw["permissions"])
    
    # 2) scope (space-delimited) or scopes (array)
    scope = raw.get("scope")
    if isinstance(scope, str):
        perms.update(s for s in scope.split() if s)
    
    # 3) Admin role -> implicit admin:all
    roles = raw.get("roles")
    if isinstance(roles, (list, tuple)):
        if any(str(r).lower() == "admin" for r in roles):
            perms.add("admin:all")
    
    return perms

def has_perms(user, any_of: Iterable[str] | str) -> bool:
    """Check if user has any of the required permissions."""
    req = _as_set(any_of)
    eff = current_permissions(user)
    if "admin:all" in eff:
        return True  # Super permission
    return any(r in eff for r in req)
\end{lstlisting}

\subsubsection{Policy File Structure}

The \texttt{roles.yaml} file defines comprehensive role policies:

\begin{lstlisting}[language=yaml,caption={Role policy definition structure}]
roles:
  admin:
    description: "Full administrative access"
    tools:
      allow: ["*"]
      deny: []
    scopes:
      - "admin:*"
      - "models:*"
      - "tools:*"
      - "db:*"
      - "graph:*"
    safety:
      output_guard: "standard"
      intent_filter: "standard"
      pii_scrubber: "redact"
    limits:
      rpm: 240
      tpm: 160000
      concurrent_jobs: 20

  researcher:
    description: "Scientists exploring the graph"
    tools:
      allow:
        - "graph.query"
        - "graph.search"
        - "catalog.discover"
      deny:
        - "graph.crud"
        - "system.*"
    scopes:
      - "models:list"
      - "models:infer"
      - "db:query:read"
    limits:
      rpm: 60
      tpm: 40000
\end{lstlisting}

\subsubsection{Per-Tool Authorization}

MCP tools declare their required scopes via the \texttt{@mcp\_tool} decorator. Authorization is checked before each tool invocation:

\begin{lstlisting}[language=Python,caption={Tool-level authorization declaration}]
@mcp_tool(
    name="graph.crud",
    description="Create, update, delete graph nodes/edges",
    required_scopes=["graph:crud", "db:write"]
)
async def graph_crud(
    operation: str,
    node_type: str,
    properties: dict,
    ctx: ToolContext
) -> ToolResult:
    # Authorization already verified by framework
    ...
\end{lstlisting}

\subsubsection{Role Inference}

For UI customization and logging, the platform can infer a user's effective role from their permissions:

\begin{lstlisting}[language=Python,caption={Role inference from permissions}]
def infer_role_from_principal(principal) -> str | None:
    """Best-effort role inference from principal scopes."""
    tokens = set()  # Collect all scopes/permissions/roles
    
    # Extract from various claim sources
    for source in [principal.raw.get("scope"), 
                   principal.raw.get("scopes"),
                   principal.raw.get("permissions")]:
        if isinstance(source, str):
            tokens |= set(source.split())
        elif isinstance(source, (list, tuple)):
            tokens |= set(str(s) for s in source)
    
    tokens = {t.lower() for t in tokens}
    
    if "admin:all" in tokens or "admin" in tokens:
        return "admin"
    if any(t in tokens for t in ("tools:all", "tools:invoke:all")):
        return "power_user"
    if any(t in tokens for t in ("tools:basic", "tools:invoke:basic")):
        return "basic"
    return None
\end{lstlisting}

\subsubsection{Authorization Error Responses}

When authorization fails, the platform returns standardized HTTP 403 Forbidden responses:

\begin{lstlisting}[language=Python,caption={Authorization denial with details}]
def authorize_or_403(user, required_scopes, **kwargs) -> None:
    """Raise HTTP 403 if authorization denies access."""
    decision = authorize(user, required_scopes, **kwargs)
    if not decision.allowed:
        raise HTTPException(
            status_code=status.HTTP_403_FORBIDDEN,
            detail=f"Not authorized: requires {list(decision.required)}",
        )
\end{lstlisting}

This design ensures that:
\begin{enumerate}
    \item Users receive clear feedback about missing permissions.
    \item All authorization decisions are audited for compliance review.
    \item The authorization logic is centralized and consistent across all endpoints.
    \item Role definitions can be updated without code changes (via YAML policies).
\end{enumerate}

